\section{Introdução}\label{sec:introduacao}

A Biometria é uma forma de identificação que se baseia nas características físicas e comportamentais das pessoas, como íris, impressões digitais e voz. Sistemas biométricos operam sobre a premissa de que tais características são únicas de cada pessoa, e, se puderem ser extraídas através de sensores, podem ser convertidas para representações numéricas e processadas por algoritmos de reconhecimento de padrões~\cite{handbook_of_fp_recognition}. Como os dados biométricos são particulares a cada pessoa eles podem ser usados como forma de autenticação e de forma mais segura do que o uso de senhas~\cite{biometrics_a_tool_for_information_security}. 

Neste projeto, será dado o enfoque na vertente física, em particular, na biometria através de impressões digitais. Essa modalidade biométrica se baseia nos padrões de pequenas elevações e depressões na pele das pontas dos dedos~\cite{handbook_of_fp_recognition}. Os usuários são então reconhecidos com base nas diferenças nesses padrões.

Conforme definido por~\citeonline{Jain2004_IntroToBiometrics}, sistemas biométricos são sistemas de reconhecimento de padrões que extraem características de dados biométricos e então comparam as características extraídas com uma referência biométrica em um banco de dados. Nesse sentido,~\citeonline{handbook_of_fp_recognition} explica que, para a autenticação através de digitais, a representação das características extraídas das imagens é um problema fundamental da verificação de digitais. 

Como não há a garantia de que as imagens da digital de um mesmo dedo serão exatamente iguais em termos de intensidade de pixel, orientação e formato, há a necessidade de se extrair características a partir das quais seja possível diferenciar digitais com uma boa acurácia e que sejam invariantes para uma dada digital. Somente assim o uso de digitais em sistemas de autenticação se tornariam viáveis. Além 

O objetivo deste trabalho é avaliar diferentes técnicas de extração de características para reconhecimento de usuários pela impressão digital. As demais seções do projeto estão organizadas da seguinte forma: na Seção~\ref{sec:fundamentacao_teorica}, são introduzidos conceitos sobre reconhecimento de usuários pela impressão digital e a revisão bibliográfica realizada; na Seção~\ref{sec:objetivos}, são definidos os objetivos; na Seção~\ref{sec:metodologia}, a metodologia é descrita, incluindo conjuntos de dados e métricas; na Seção~\ref{sec:resultados} são exibidos os resultados obtidos e na Seção~\ref{sec:conclusao} é feita a conclusão do relatório.
